\section{Reproducing Example Calculations}
\label{sec:successful-examples}

The example data catalogue maps each figure to its numerical inputs and
plotting script. The corresponding commands are collected in
\path{analysis/README.md}. The supplied caches contain physical profiles, particle projections,
moments, and coordinate clocks. Plotting reads these arrays directly, so a
figure can be regenerated without repeating the simulation.
The movie catalogue in \path{movies/README.md} describes each physical
evolution and defines the plotted observables.

\begin{center}
\small
\begin{tabular}{ll}
\hline
Physical calculation & Saved numerical data \\
\hline
Expanding subhalo VP/BP & \path{movies/data/halos}, \path{analysis/data/phase_through10000} \\
Physical virial moments & \path{analysis/data/virial_summary} \\
Cold collisionless collapse & \path{analysis/data/cold_vp} \\
Static and expanding Gaussian VP/BP & \path{analysis/data/gaussian_vp} \\
Material SIDM evolution & \path{analysis/data/sidm} \\
Relativistic shell collapse & \path{analysis/data/relativistic} \\
\hline
\end{tabular}
\end{center}

The Gaussian generator in
\path{cpu/tools/generate_bound_gaussian_shell_restart.py} writes the same
physical particle realization in static and expanding coordinates. Its
companion example in
\path{cpu/examples/modular/vp_bound_gaussian_shell_gpu.md} explains the
initial distribution, physical normalization and solver commands.

For another physical problem, choose the equation and initial distribution
first. Then set the boundaries, material law, and numerical resolution in the
corresponding solver interface.
The input-format examples under \path{cpu/examples/modular} and the custom
sampling script \path{cpu/examples/user_distribution.py} explain the grid
and particle conventions. The initial-condition chapter derives their
normalization and the model-specific restrictions.
